% !TEX root = ../paper.tex

\section{Related Work}

\subsection{Classical Image Registration}

Classical image registration pipelines typically rely on hand-crafted local features followed by
robust geometric fitting. Representative methods built on SIFT, SURF, ORB, BRISK, KAZE, or AKAZE
remain attractive because they are simple, interpretable, and do not require task-specific
training
\cite{lowe2004sift,bay2008surf,rublee2011orb,leutenegger2011brisk,alcantarilla2012kaze,alcantarilla2013akaze}.
However, their performance can degrade substantially in cross-modal settings, especially when the
appearance gap between visible and thermal imagery reduces descriptor repeatability or when thermal
frames contain weak texture and large homogeneous regions.

These methods remain important baselines in our study because they represent the classical
registration view of the problem: if sufficiently stable local features can be extracted, a single
pairwise homography may already solve the task. Robust model estimation remains central in these
pipelines, from RANSAC and multiple-view geometry to more recent variants such as MAGSAC and
MAGSAC++
\cite{fischler1981ransac,hartley2004multiple,barath2019magsac,barath2020magsacpp}. Our results
suggest that the single-pair assumption is too weak for real UAV RGB-thermal scenes, where
scene-level reliability often remains difficult even when some individual pairs are usable.

\subsection{Deep Pairwise Matching for Multimodal Registration}

Recent dense and semi-dense matching methods have significantly improved pairwise geometric
correspondence estimation. Learned local features have evolved from trainable detectors and
descriptors such as SuperPoint, LF-Net, D2-Net, R2D2, DISK, S2DNet, HardNet, L2-Net, SOSNet,
GeoDesc, ContextDesc, ASLFeat, ALIKE, and Key.Net
\cite{detone2018superpoint,ono2018lfnet,dusmanu2019d2net,revaud2019r2d2,tyszkiewicz2020disk,germain2020s2dnet,mishchuk2017hardnet,tian2017l2net,tian2019sosnet,luo2018geodesc,luo2019contextdesc,luo2020aslfeat,zhao2023alike,barroso2019keynet},
while match filtering and end-to-end correspondence reasoning have progressed through
SuperGlue, NCNet, Patch2Pix, COTR, LoFTR, ASpanFormer, LightGlue, DKM, RoMa, Efficient LoFTR,
OmniGlue, MINIMA, and XoFTR
\cite{sarlin2020superglue,rocco2018ncnet,zhou2021patch2pix,jiang2021cotr,sun2021loftr,chen2022aspanformer,lindenberger2023lightglue,edstedt2023dkm,edstedt2024roma,wang2024efficientloftr,jiang2024omniglue,ren2025minima,tuzcuoglu2024xoftr}.
In particular, these methods are better at recovering useful correspondences in weak-texture
regions, repetitive structures, or substantial appearance shifts. Broad reviews of this transition
from classical pipelines to trainable matching, together with practical benchmarking studies,
provide additional context
\cite{ma2021survey,liao2024survey}.

Our work builds directly on this progress instead of trying to replace it. The empirical picture in
UAV-TAlign-1K is that strong off-the-shelf multimodal matchers already provide very high pairwise
availability. This changes the central question. The bottleneck is no longer only whether a matcher
can produce correspondences for a given RGB-thermal pair, but whether those frame-level estimates
can be converted into a stable and reliable scene-level alignment under realistic UAV operating
conditions.

\subsection{Visible-Thermal Registration and Evaluation}

Visible-thermal registration has been studied under a range of assumptions, from same-view or
near-same-view alignment to broader multimodal correspondence learning. Radiation-variation
insensitive matching methods such as RIFT and deep alignment or stitching formulations such as
CLKN and UDIS indicate the breadth of existing registration formulations
\cite{li2020rift,chang2017clkn,nie2021udis}. Many existing settings either operate under relatively
controlled viewpoint changes or evaluate pairwise registration in a way that does not require a
scene-level acceptance decision. By contrast, practical UAV RGB-thermal data frequently combine
cross-FOV capture, cross-resolution sensing, thermal rendering variation, and low-light conditions.
These factors make scene-level reliability a more central concern than in many standard
visible-thermal evaluation setups.

A further complication is that point-level manual geometric ground truth is expensive and sometimes
ill-posed in weak-texture thermal frames. For the current PRCV scope, we therefore avoid presenting
our dataset package as a fully supervised benchmark. Instead, we package a public subset,
UAV-TAlign-1K, together with a practical label-free evaluation protocol that measures
scene-consistency and relative alignment quality without requiring manual landmark annotation as a
prerequisite.

\subsection{UAV Multimodal Datasets and the Remaining Gap}

UAV multimodal datasets have supported progress in detection, tracking, segmentation, and broader
cross-modal perception, spanning aligned pedestrian and automotive benchmarks as well as
RGB-thermal detection and tracking benchmarks such as KAIST, FLIR, LLVIP, M3FD, DroneVehicle,
VTUAV, LasHeR, and Anti-UAV
\cite{hwang2015kaist,flir2018thermal,jia2021llvip,liu2022tardal,sun2022dronevehicle,zhang2022vtuav,li2022lasher,jiang2023antiuav}.
However, registration is not always the primary task in those benchmarks, and the protocol
assumptions needed for scene-level RGB-thermal alignment are often different from those needed for
downstream semantic tasks. In particular, a practical registration pipeline must decide not only
how to align frames, but also when the scene-level evidence is strong enough to trust the
resulting transform.

This is the gap our paper targets. Rather than introducing a new matcher architecture, we focus on
a practical scene-level registration setting for real UAV RGB-thermal imagery. Our contribution
lies in the combination of a reliability-oriented pipeline, a public subset packaging, and a
label-free protocol that makes the problem studyable without first building dense manual geometric
annotations.
